\section{Introduction}

The Efficient Market Hypothesis posits that asset prices instantaneously incorporate all available information, causing price changes to follow an unpredictable martingale process \cite{fama1970efficient}. However, empirical financial returns systematically display stylized non-equilibrium anomalies, including heavy-tailed distributions, volatility clustering, and macroscopic herding cascades \cite{lux1999scaling, cont2001empirical, shiller2000irrational}. These observations indicate that markets operate as complex open systems where boundedly rational agents process historical signals over finite cognitive horizons and interact through social observation networks \cite{bouchaud2013more, farmer2009economy, kirman1993ants}.

A fundamental challenge in quantitative portfolio management is backtest overfitting. Strategies optimized by in-sample parameter selection frequently collapse out-of-sample \cite{bailey2017probability}. Traditional quantitative finance models backtest overfitting as a purely statistical error arising from multiple hypothesis testing without modeling the physical feedback dynamics that cause strategies to lock into stale market regimes.

Recent developments in non-equilibrium statistical physics have shown that swarms of decision-making agents equipped with finite First-In-First-Out (FIFO) informational queues undergo macroscopic phase transitions \cite{li2026informational}. Beyond a critical memory depth $M_c$, interaction dynamics cause symmetric populations to spontaneously break symmetry, locking the crowd into persistent consensus states.

This paper establishes the Physics-Informed Market Strategy (PIMS) framework by bridging micro-level quantitative strategy backtesting with macro-level swarm physics on the Indian Nifty 50 Total Return Index (TRI) from 2015 to 2025. The core theoretical discovery is that quantitative strategy overfitting ($\Delta \text{Sharpe}$) is the direct empirical manifestation of a macroscopic supercritical pitchfork bifurcation ($k < 0, M > M_c \approx 9.82$ trading days) in the underlying crowd dynamics.

\subsection{Core Contributions}
\begin{enumerate}
    \item \textbf{Microscopic Sentiment Injection vs. Macro Aggregation}: Evaluating five behavioral memory kernels (Pushover, Opportunist, Traditionalist, Contrarian, and Curmudgeon) across 852 configurations, we show that queue-level sentiment injection achieves out-of-sample Sharpe ratios up to 1.486 (annualized return +18.4\%, maximum drawdown -14.6\%), significantly outperforming linear blending (Sharpe 0.627) and majority voting (Sharpe 0.153). Furthermore, injecting Contrarian signals into a Pushover host achieves Sharpe 1.470 with daily turnover reduced to 15.2\%.
    \item \textbf{Phase Transitions in Empirical Crowds}: Calibrated on daily Nifty 50 returns, agent swarms ($N=150$) undergoing peer interaction ($p_{\text{peer}}=1.0$) cross the pitchfork bifurcation boundary between $M=9$ and $M=11$ days, matching the analytical kinetic prediction $M_c = \frac{\pi}{2(1-2\eta)^2} \approx 9.82$ days for empirical noise $\eta = 0.30$. Conversely, isolated price observation ($p_{\text{peer}}=0$) maintains strictly positive potential curvature ($k > 0$), proving that social contagion is a necessary condition for herding.
    \item \textbf{Critical Slowing Down and Circuit Breaking}: In panic relaxation simulations initialized from unanimous bearish consensus ($n_+(0)=0$), relaxation half-lives $T_{1/2}$ diverge near $M_c$. Introducing a minority fraction ($\phi_C \ge 10\%$) of structural macro trend anchors breaks potential symmetry, guaranteeing panic recovery within 30 to 60 trading days.
    \item \textbf{Physics-Informed Strategy Selection}: Enforcing the physical sub-critical constraint ($M_{\text{eff}} \le M_c$) as a pre-fit prior filter increases mean out-of-sample Sharpe ratios by +41.4\% (0.564 versus 0.399) and reduces the overfitting gap by 37.7\% (0.375 versus 0.602) relative to naive in-sample selection.
\end{enumerate}
